% Partie : counterfactuals
% Chapitre : introduction
% Section : material-conditional

\documentclass[../../../include/open-logic-section]{subfiles}

\begin{document}

\olfileid[fr]{cnt}{int}{mat}

\olsection{Le conditionnel matériel}

En anglais, la forme la plus simple d'un énoncé conditionnel est « If
\dots then \dots » (en français, « Si \dots, alors \dots ») : les
pointillés représentent eux-mêmes des énoncés. Par exemple : « Si le
majordome est le coupable, alors le jardinier est innocent. » Dans un
cours élémentaire de logique, on apprend à symboliser ces
conditionnels au moyen du connecteur~$\lif$ : les parties indiquées
par les pointillés sont représentées, par exemple, par les
!!{formula}s $!A$ et~$!B$, et le conditionnel entier par
$!A \lif !B$.

Le connecteur~$\lif$ est \emph{vérifonctionnel} : la valeur de
vérité---$\True$ ou $\False$---de $!A \lif !B$ est déterminée par
les valeurs de vérité de $!A$ et de~$!B$. La formule $!A \lif !B$
est vraie si et seulement si $!A$ est fausse ou $!B$ est vraie ; elle
est fausse dans le seul autre cas. Pour une valuation~$\pAssign v$,
nous définissons $\pSat{v}{!A \lif !B}$ si et seulement si
$\pSat/{v}{!A}$ ou $\pSat{v}{!B}$. Le connecteur~$\lif$ ainsi
interprété est appelé \emph{conditionnel matériel}.

Cette définition entraîne plusieurs faits logiques élémentaires.
D'abord, le théorème de déduction vaut pour le conditionnel matériel :
\begin{align}
  \text{Si } \Gamma, !A \Entails !B \text{ alors } \Gamma \Entails !A \lif !B
\end{align}
Il est vérifonctionnel : $!A \lif !B$ et $\lnot !A \lor !B$ sont
équivalents :
\begin{align}
  !A \lif !B & \Entails \lnot !A \lor !B\\
  \lnot !A \lor !B & \Entails !A \lif !B
  \intertext{Un conditionnel matériel découle de son conséquent et de
    la négation de son antécédent :}
  !B & \Entails !A \lif !B\\
  \lnot !A & \Entails !A \lif !B
  \intertext{La négation d'un conditionnel matériel équivaut à la
    conjonction de son antécédent et de la négation de son conséquent :
    si $!A \lif !B$ est faux, $!A \land \lnot !B$ est vrai, et
    réciproquement :}
  \lnot(!A \lif !B) & \Entails !A \land \lnot !B\\
  !A \land \lnot !B & \Entails \lnot(!A \lif !B)
  \intertext{Le conditionnel matériel admet le modus ponens :}
  !A, !A \lif !B & \Entails !B
  \intertext{Il permet de regrouper les conséquents :}
  !A \lif !B,  !A \lif !C & \Entails !A \lif (!B \land !C)
  \intertext{On peut toujours renforcer l'antécédent (propriété dite
    ici de \emph{monotonie}) :}
  !A \lif !B & \Entails (!A \land !C) \lif !B
  \intertext{Le conditionnel matériel est transitif : la règle
    d'enchaînement est valide :}
  !A \lif !B, !B \lif !C & \Entails !A \lif !C
  \intertext{Il équivaut à sa contraposée :}
  !A \lif !B & \Entails \lnot !B \lif \lnot !A\\
  \lnot !B \lif \lnot !A & \Entails !A \lif !B
\end{align}

Toutes ces inférences sont utiles et ne posent pas de difficulté dans
le raisonnement mathématique. Dans la littérature philosophique et
linguistique, on trouve toutefois de nombreux contre-exemples
allégués aux inférences correspondantes dans des contextes non
mathématiques. Ils suggèrent que le conditionnel matériel~$\lif$
n'est pas---ou du moins pas toujours---le connecteur approprié pour
symboliser les énoncés anglais de la forme « If \dots then \dots ».

\end{document}
